%==============================================================
%  gbm.tex
%  hasklib — mathematical companion / derivations
%
%  Geometric Brownian Motion: exact terminal law via Ito's lemma
%  applied to ln(X_t), and the moments used by the moment-matching
%  test in test_stochastic.cpp.
%
%  Depends on ito_foundations.tex (Ito's lemma, Theorem 1).
%  Downstream of nothing; ou.tex and cir.tex are siblings, not
%  dependents, of this file.
%==============================================================
\documentclass[11pt]{article}

%---------------- packages ----------------
\usepackage[utf8]{inputenc}
\usepackage[T1]{fontenc}
\usepackage{amsmath, amssymb, amsthm}
\usepackage{mathtools}
\usepackage[a4paper, margin=1in]{geometry}
\usepackage{xcolor}
\usepackage{hyperref}
\usepackage{cleveref}
\usepackage{bm}
\usepackage{harvard}   % Harvard author-date referencing (with agsm.bst)

%---------------- notation macros ----------------
% (kept identical to ito_foundations.tex so the two files read as one document)
\newcommand{\Prob}{\mathbb{P}}
\newcommand{\Qmeas}{\mathbb{Q}}
\newcommand{\E}{\mathbb{E}}
\newcommand{\Var}{\operatorname{Var}}
\newcommand{\Cov}{\operatorname{Cov}}
\newcommand{\filt}{\mathcal{F}}
\newcommand{\Om}{\Omega}

\newcommand{\dd}{\mathrm{d}}
\newcommand{\dW}{\dd W_t}
\newcommand{\dt}{\dd t}
\newcommand{\R}{\mathbb{R}}

\newcommand{\drift}{a}
\newcommand{\diff}{b}

%---------------- theorem environments ----------------
\theoremstyle{plain}
\newtheorem{theorem}{Theorem}
\newtheorem{lemma}{Lemma}
\newtheorem{proposition}{Proposition}
\newtheorem{corollary}{Corollary}

\theoremstyle{definition}
\newtheorem{definition}{Definition}

\theoremstyle{remark}
\newtheorem{remark}{Remark}

\newcommand{\TODO}[1]{\medskip\noindent\textbf{[\textcolor{red}{TO DERIVE}: #1]}\medskip}

%==============================================================
\title{Geometric Brownian Motion: Exact Terminal Sampler\\
       \large \texttt{hasklib} mathematical companion}
\author{Hubert}
\date{\today}
%==============================================================

\begin{document}
\maketitle

\begin{abstract}
This note derives the result that \texttt{GBM::sample\_terminal} relies
on: applying It\^o's lemma to $\ln X_t$ linearises the geometric
Brownian motion SDE and yields an exact (not discretised) sampler for
$X_T$ given $X_0$. We also derive the first two moments, which are
what \texttt{test\_stochastic.cpp} checks against simulated output.
This file assumes It\^o's lemma as established in
\texttt{ito\_foundations.tex} (Theorem~1 there); it is not re-derived
here. The derivation is classical; see \citeasnoun{shreve2004} for the
standard textbook treatment.
\end{abstract}

%--------------------------------------------------------------
\section{The GBM SDE}
%--------------------------------------------------------------

\begin{definition}[Geometric Brownian motion]\label{def:gbm}
$X_t$ follows geometric Brownian motion with drift $\mu\in\R$ and
volatility $\sigma>0$ if it solves
\begin{equation}\label{eq:gbmsde}
   \dd X_t = \mu X_t\,\dt + \sigma X_t\,\dW, \qquad X_0>0.
\end{equation}
In the notation of \texttt{ito\_foundations.tex} this is the choice
$\drift(t,x)=\mu x$, $\diff(t,x)=\sigma x$ in the general SDE
$\dd X_t = \drift(t,X_t)\dt+\diff(t,X_t)\dW$.
\end{definition}

\begin{remark}[Correspondence with the code]
This is exactly the pair of functions \texttt{GBM::drift} and
\texttt{GBM::diffusion} return: \texttt{drift(t,x)} $=\mu x$,
\texttt{diffusion(t,x)} $=\sigma x$. \texttt{EulerMaruyama} consumes
these polymorphically through the \texttt{StochasticProcess}
interface; the sampler derived below is the exact alternative
to that discretisation, used by \texttt{GBM::sample\_terminal} instead
of stepping the Euler scheme forward.
\end{remark}

Both $\drift$ and $\diff$ are linear in $x$ and (for $x>0$, which GBM
never leaves since $X_t=X_0\exp(\cdot)>0$ once we derive the solution
below) globally Lipschitz on any bounded interval away from $0$, so the
existence--uniqueness assumption of \texttt{ito\_foundations.tex}
applies and \eqref{eq:gbmsde} has a unique strong solution.

%--------------------------------------------------------------
\section{Linearising via It\^o's lemma}
%--------------------------------------------------------------

The SDE~\eqref{eq:gbmsde} has no constant-coefficient closed form as
written, because both coefficients scale with $X_t$. The standard fix
is to apply It\^o's lemma to $f(t,x)=\ln x$: this removes the
$x$-dependence from both coefficients and turns \eqref{eq:gbmsde} into
an SDE with constant drift and diffusion, which integrates
directly.

\begin{theorem}[Exact terminal law of GBM]\label{thm:gbm-exact}
Let $X_t$ solve \eqref{eq:gbmsde}. Then for any $T>0$,
\begin{equation}\label{eq:gbmexact}
   X_T = X_0\exp\!\left(\Big(\mu-\tfrac12\sigma^2\Big)T + \sigma W_T\right),
\end{equation}
where $W_T\sim\mathcal N(0,T)$.
\end{theorem}

\begin{proof}
Apply It\^o's lemma (Theorem~1 of \texttt{ito\_foundations.tex}) to
$\varphi(t,x)=\ln x$, so $\frac{\partial\varphi}{\partial t}=0$, $\frac{\partial\varphi}{\partial x}=\frac{1}{x}$, $\frac{\partial^2\varphi}{\partial x^2}=-\frac{1}{x^2}$. With
$\drift(t,X_t)=\mu X_t$ and $\diff(t,X_t)=\sigma X_t$,
\[
   \dd(\ln X_t)
     = \Big( 0 + \mu X_t\cdot\frac1{X_t}
             + \tfrac12\,\sigma^2 X_t^2\cdot\Big(\!-\frac1{X_t^2}\Big) \Big)\dt
       + \sigma X_t\cdot\frac1{X_t}\,\dW .
\]
The $X_t$ factors cancel completely, leaving
\begin{equation}\label{eq:lnlinear}
   \dd(\ln X_t) = \Big(\mu-\tfrac12\sigma^2\Big)\dt + \sigma\,\dW .
\end{equation}
The right-hand side no longer depends on $X_t$: it is Brownian motion
with constant drift $\mu-\tfrac12\sigma^2$ and constant volatility
$\sigma$, so \eqref{eq:lnlinear} integrates directly from $0$ to $T$:
\[
   \ln X_T - \ln X_0
     = \Big(\mu-\tfrac12\sigma^2\Big)T + \sigma\,(W_T-W_0)
     = \Big(\mu-\tfrac12\sigma^2\Big)T + \sigma W_T,
\]
using $W_0=0$. Exponentiating gives \eqref{eq:gbmexact}. Since $W_T\sim
\mathcal N(0,T)$ by definition of Brownian motion, this is a complete
description of the law of $X_T$ given $X_0$.
\end{proof}

\begin{remark}
Equation~\eqref{eq:gbmexact} is sampled with a single normal draw: generate
$Z\sim\mathcal N(0,1)$, set $W_T=\sqrt T\,Z$, and evaluate the
right-hand side. No time-stepping is needed, which is what makes it
an exact sampler rather than a discretisation, and it is exactly
what \texttt{GBM::sample\_terminal} implements using
\texttt{NormalRng::draw()}.
\end{remark}

\begin{remark}[Where the $-\tfrac12\sigma^2$ comes from]
The drift that appears in \eqref{eq:gbmexact} is not $\mu$, the
drift of $X_t$ itself, but $\mu-\tfrac12\sigma^2$. That correction is
exactly the $\tfrac12\diff^2 \frac{\partial^2\varphi}{\partial x^2}$ term of It\^o's lemma, which is the term
with no analogue in ordinary calculus, discussed at length in
\texttt{ito\_foundations.tex}. If one (incorrectly) applied the
ordinary chain rule to $\ln X_t$, the correction would be missing and
the resulting sampler would be biased upward in expectation, as
\cref{prop:gbm-mean} below makes precise.
\end{remark}

%--------------------------------------------------------------
\section{Moments}
%--------------------------------------------------------------

The moment-matching test in \texttt{test\_stochastic.cpp} checks
simulated terminal samples against the following closed-form mean and
variance.

\begin{lemma}[Moment generating function of a Gaussian]\label{lem:mgf}
\leavevmode\\
If $Y\sim\mathcal N(m,v)$ then $\E[e^Y]=\exp\!\big(m+\tfrac12 v\big)$.
\end{lemma}


\begin{proof}
\leavevmode\\
Write $Y=m+\sqrt v\,Z$ with $Z\sim\mathcal N(0,1)$. Complete the square
in the Gaussian integral:
\[
   \E[e^Y] = e^m\,\E\big[e^{\sqrt v\,Z}\big]
           = e^m \int_{-\infty}^{\infty}
               \frac{1}{\sqrt{2\pi}}\,e^{\sqrt v\,z-\frac12 z^2}\,\dd z
           = e^m\,e^{\frac12 v}
             \int_{-\infty}^{\infty}
               \frac{1}{\sqrt{2\pi}}\,e^{-\frac12(z-\sqrt v)^2}\,\dd z
           = e^{m+\frac12 v},
\]
the last integral being $1$ since it is a $\mathcal N(\sqrt v,1)$
density integrated over $\R$.
\end{proof}

\begin{proposition}[Mean of GBM]\label{prop:gbm-mean}
$\E[X_T] = X_0\,e^{\mu T}$.
\end{proposition}

\begin{proof}
\leavevmode\\
By \cref{thm:gbm-exact}, $X_T=X_0\,e^Y$ with
$Y=(\mu-\tfrac12\sigma^2)T+\sigma W_T$. Since $W_T\sim\mathcal N(0,T)$,
$Y\sim\mathcal N\big((\mu-\tfrac12\sigma^2)T,\ \sigma^2 T\big)$. Applying
\cref{lem:mgf} with $m=(\mu-\tfrac12\sigma^2)T$ and $v=\sigma^2 T$,
\[
   \E[X_T] = X_0\,\E[e^Y]
           = X_0\exp\!\Big( \big(\mu-\tfrac12\sigma^2\big)T + \tfrac12\sigma^2 T \Big)
           = X_0\,e^{\mu T}.
\]
The $-\tfrac12\sigma^2 T$ It\^o correction and the $+\tfrac12\sigma^2 T$
from the moment generating function cancel exactly, leaving the mean
growing at the "natural" rate $\mu$ despite the correction appearing
inside every individual path.
\end{proof}

\begin{proposition}[Variance of GBM]\label{prop:gbm-var}
$\Var(X_T) = X_0^2\,e^{2\mu T}\big(e^{\sigma^2 T}-1\big)$.
\end{proposition}

\begin{proof}
$\Var(X_T)=\E[X_T^2]-\big(\E[X_T]\big)^2$. For the second moment, write
$X_T^2=X_0^2\,e^{2Y}$ with $Y$ as above, so $2Y\sim\mathcal
N\big(2(\mu-\tfrac12\sigma^2)T,\ 4\sigma^2 T\big)$. By \cref{lem:mgf}
with $m=2(\mu-\tfrac12\sigma^2)T$ and $v=4\sigma^2 T$,
\[
   \E[X_T^2] = X_0^2\exp\!\Big( 2\big(\mu-\tfrac12\sigma^2\big)T + 2\sigma^2 T \Big)
             = X_0^2\,e^{(2\mu+\sigma^2)T}.
\]
Subtracting $\big(\E[X_T]\big)^2=X_0^2\,e^{2\mu T}$ from
\cref{prop:gbm-mean} and factoring,
\[
   \Var(X_T) = X_0^2\,e^{(2\mu+\sigma^2)T} - X_0^2\,e^{2\mu T}
             = X_0^2\,e^{2\mu T}\big(e^{\sigma^2 T}-1\big). \qedhere
\]
\end{proof}

\begin{remark}[What the test checks]
\texttt{test\_stochastic.cpp} draws a large number of terminal samples
from \texttt{GBM::sample\_terminal}, computes their sample mean and
sample variance, and checks these against
\cref{prop:gbm-mean,prop:gbm-var} to within Monte Carlo tolerance. This
is a check on the sampler, not on \eqref{eq:gbmsde} itself ---
it would catch, for instance, a sign error in the $-\tfrac12\sigma^2$
correction, since that error changes the mean but leaves the
qualitative shape of the distribution intact.
\end{remark}

%--------------------------------------------------------------
\bibliographystyle{agsm}
\bibliography{refs}

\end{document}
